\section{Limitations and Threats to Validity}
\label{sec:limitations}
% ================================================================================

\paragraph{Boundaries of applicability.} The Bridge Contract adds value
only when the client can detach or the human can be unavailable. In
scenarios with perfect-availability connectivity and an always-available
human, the Bridge Contract is overhead: every state transition emits an
event, and every detach produces a resume record that is
never used. Conformant deployments MAY skip resume-record production when
no detach is possible, but MUST still emit lifecycle events
(ATTACHED, TERMINATION).

\paragraph{What the bridge does not see.} The contract covers what
crosses the bridge (Section~\ref{sec:coverage}). Where the bridge sees
an approval request but not the tool execution, I2 applies to the
request and execution is outside the contract. A runtime that runs a
tool without first submitting an approval request is outside the gate in
the same way: I2 and I13 govern only the calls that cross the bridge.
Where the runtime
announces its own authority or mode change, I3 is \emph{observed}, not
enforced: the bridge reports the change and SHOULD instruct a revert,
but the runtime holds the widened authority until it complies. A bridge
claims \emph{enforced} authority only where the runtime routes all
authority changes through it; otherwise its authority is
\emph{observed} (Section~\ref{sec:coverage}). Neither
limit can be removed by a bridge alone; both need a runtime that routes
execution and authority changes through it.

\paragraph{Post-quantum signatures.} The contract fixes what a key must
prove at re-attach (possession, bound to the session, the nonce and the
acknowledged prefix of the chain), not which signature scheme does the
proving. The implementations and the evidence use Ed25519 and, as a second
scheme, ECDSA P-256 (V59--V62); neither is post-quantum secure, and the model checking treats a signature as a
possession predicate whose unforgeability is assumed. A re-attach signature matters only while its nonce is live
($T_{\text{nonce}}$, default 60~s), so a later break of Ed25519 does
not reopen past re-attaches; but a break that is feasible while a
session is still live would let an attacker who knows a registered
public key forge a re-attach for any session that is still in a
detached or re-attaching cell, which in the \emph{persistent} profile
is up to $T_{\text{total}}$ of that profile (Table~\ref{tab:timer-profiles}) after the last detach. The signatures
that must stay unforgeable for years are those on archived audit evidence,
the operator signatures on events and checkpoints (I-SIGN-EVENT,
I-CHECKPOINT); a deployment that keeps such evidence for long should move
those to a post-quantum or hybrid scheme first. Section~\ref{sec:ext-schemes}
gives the scheme-list convention and the guidance as an extension that
weakens no MUST invariant; no post-quantum scheme is implemented or tested.

\paragraph{Out of scope.} The following are explicitly out of scope
(see \S\ref{sec:scope} for the full failure model): bridge
replication and failover, network partition (partial detach), true
concurrency on the same session, compromise of the client or bridge
keys, Byzantine bridge, multi-hop composition (BC $\circ$ BC), and
multi-agent composition (A2A, MCP). Auditor
marking (annotations anchored to events via \texttt{event\_hash}),
routing, and broadcast are also explicitly out of scope for this
version; BC reserves \texttt{AUDITOR\_MARK}, \texttt{ROUTING\_HINT},
and \texttt{BROADCAST\_SUBSCRIPTION} as extension points for a
separate specification. A
future revision may extend the contract to cover some of these.

\paragraph{Implementation divergence.} Two bridge implementations may
satisfy the invariants but produce wire formats that are not
byte-for-byte interoperable. The Bridge Contract defines the resume record
(Definition~\ref{def:rt}) as a tuple, the supplementary schemas fix
the envelope and the payload of every event kind
(Section~\ref{sec:conformance}), and Appendix~\ref{app:wire-format}
specifies the client--bridge wire format (the op-types,
single-envelope JSON, signature scope, error reasons); it is normative
for the signed bytes and reason names, which the vectors and known
answers check, and for the envelope only where the wire front-ends of the
two implementations serve it. What remains
open is the transport binding (HTTP, WebSocket, QUIC): the wire
format is the JSON message content, and the transport carries it.
The two implementations in this deposit agree on the committed event
kinds and the final cell on every vector V01--V134 that does not require
checkpoint emission (V52--V54, V88, V89). In wire mode
(\texttt{run\_vectors.py -{}-wire}) the vector actions that have a
server-side handler in the wire front-end of each implementation, which
includes session creation (\texttt{attach}), \texttt{resolve}, \texttt{extend} and the
operator events, are sent as
Appendix~B messages with real signatures and the bridge verifies them from the
envelope; the actions without a handler (clock and restart hooks, log and
registry tampering, a re-attach left incomplete or flagged as a transport
replay) are applied to the bridge directly (the Python runner also checks that they encode and decode).
The byte level is checked separately: the
canonical encoding is RFC~8785 for strings, integers up to $2^{53}$ and
finite numbers, and the event hash, the content hash and the
re-attach and join proofs are pinned by known-answer vectors
(\texttt{vectors/bytes\_kat.json}, T101--T104) that both bridges read.
Outside that subset (NaN, infinities, integers beyond $2^{53}$, which
no event of the contract carries) a bridge refuses to canonicalise.
Byte-level agreement with a third implementation has not been shown,
only with the JavaScript string and number rules and the Node.js Ed25519
verifier on the file's contents.

\paragraph{The human channel.} The contract does not say who may emit
\texttt{HUMAN\_RESOLVED}, \texttt{TOOL\_CALL\_ANSWERED} or a clean end of
session, nor under which key. The human channel is outside the contract:
the bridge trusts its runtime and the deployment to deliver them from the
human. I14 binds an answer to the bytes the client reports having shown;
it protects against rendering drift by an honest client, not against a
client that reports whatever hash it was given.

\paragraph{Confidentiality of the log and the replay.} The resume
record and the events replayed on re-attach carry the authority scope
and the pending tool calls. In regulated deployments these may leak
sensitive information. Operators SHOULD encrypt the log at rest with a
key bound to the principal's identity, and the replay MUST travel over
an authenticated, confidential transport. The event log and the replay
are, as a deployment requirement, readable only by parties the
deployment has authorised (the principal, auditors with a documented
scope, and the bridge operator), and the deployment documents who those
parties are. The deployment enforces a retention bound on the event log
of every session, and the operator defines the policy by which the log
is archived or destroyed after it. These are deployment requirements
(Table~\ref{tab:deployment-requirements}); none is an invariant of the contract and none
is checked by it. A future revision may specify
a standard encryption envelope.

\paragraph{Audit levels.} The audit guarantees of this contract
fall into the levels below, in increasing strength:
\begin{description}[leftmargin=2em]
\item[Crash durability (I12).] The durable append of I12 and the
recovery rules of Section~\ref{sec:recovery} ensure that a bridge
crash loses no committed event and that the cell recovered from the
log matches the cell the bridge last acknowledged. This is the first
level, and the baseline MUSTs require it.
\item[Internal consistency (I5).] The hash chain of I5 makes the log
tamper-evident against an unsophisticated editor for edits inside the
log: an insertion, a deletion, a rewrite or a reordering of records
that are followed by other records breaks the chain and is detected on
verification, whether at re-attach (I4) or at restart recovery
(Section~\ref{sec:recovery}). It does not detect truncation or the
rewriting of the last events: a log cut back to an earlier event is a
valid chain, and recovery then accepts it, so a terminated session can
become active again and a revoked key can be accepted again. Closing
that needs a head (sequence and hash) held outside the log, such as an
I-CHECKPOINT anchor; the reference bridges keep none. This is also
required by the baseline MUSTs.
\item[External tamper evidence (I-SIGN-EVENT, I-CHECKPOINT).] A
determined operator with write access to the log can rewrite it and
recompute the chain. Detecting that requires operator-signed events
(I-SIGN-EVENT, SHOULD) plus checkpoints anchored outside the
operator's control (I-CHECKPOINT, SHOULD), for example in a SCITT
transparency log. This external-tamper-evidence level is not required by the baseline
MUSTs; a deployment that needs it adopts both SHOULDs.
\end{description}
The baseline contract delivers crash durability and internal
consistency; external tamper evidence is the SHOULD-level extension
that defends against a Byzantine operator (A4 is otherwise assumed).

\paragraph{Performance considerations.} Verification of the event hash
chain on re-attach is $O(n)$ in the number of events emitted since
session attach. For long-running sessions ($n > 10^4$), the bridge
SHOULD emit signed Merkle checkpoints (I-CHECKPOINT,
Appendix~\ref{app:merkle}) so that verification covers only the suffix
after the latest verified checkpoint; the conformance suite still uses
an $O(n)$ full-chain verification. The size of the event log is unbounded in this document;
the operator defines the retention policy, as a deployment requirement (above). The size of the pending
tool-call queue $Q$ in the resume record is bounded by $Q_{\max}$ (I13);
a call arriving at a full queue is refused rather than evicting an
older one.

\paragraph{Non-goals.} The contract is not a transport or message
format and does not replace AGTP, GAR, or PACE; the client--bridge
wire format is specified normatively in
Appendix~\ref{app:wire-format}, but the transport binding (HTTP,
WebSocket, QUIC) is not; the contract does not defend against a compromised bridge, a compromised
runtime, or a stolen client key; it gives no guarantee for work the
runtime performs without submitting it to the bridge (Section~\ref{sec:coverage});
and it is not a proof of correctness of any implementation.

\paragraph{Limitations register.}
Table~\ref{tab:limitations} lists each limitation of this work, what
it affects, who owns the residual risk, and the mitigation or future
work that addresses it. The register is the single place to look for
``what this document does not claim''; the prose paragraphs above
give the reasoning, the table gives the index.

\footnotesize
\begin{longtable}{p{0.18\columnwidth} p{0.22\columnwidth} p{0.18\columnwidth} p{0.26\columnwidth}}
\caption{Limitations register.}
\label{tab:limitations}\\
\toprule
\textbf{Limitation} & \textbf{Affects} & \textbf{Owner} & \textbf{Mitigation} \\
\midrule
\endfirsthead
\toprule
\textbf{Limitation} & \textbf{Affects} & \textbf{Owner} & \textbf{Mitigation} \\
\midrule
\endhead
Single author, two implementations &
Validity of the evidence (consistency, not independent confirmation) &
Author; future implementers &
External implementer; future work; V01--V134 are the interface \\
Bounded TLC (two-session run does not fire Disconnect/Resolve/Submit/Drain) &
Generalisability of the model check to multi-session detach paths &
Author; future model-checking work &
Deeper single-session runs cover the detach paths; unbounded argument is future work \\
Cryptographic checks abstracted in TLA+ (I4, I5, I14 hold only as flag abstractions; I11 is checked against an attacker that holds no client key) &
Formal rigour: unforgeability and hash-chain integrity are assumed, not model-checked; \texttt{NoForgery} checks that an accepted signature was produced by an honest client, not that the request was presented by one, so it admits a replayed honest signature; the state invariant for I6 is a tautology and the no-eviction property of I11 (f) reads over an empty set in every configuration (\texttt{LiveLimit} is one), so these two are not meaningfully TLC-checked &
Author; formal-methods community &
Conformance suite and mutants exercise the cryptography (Ed25519, SHA-256); a property that every accepted request is honest, a real I6 conjunct and a configuration with more than one live nonce are planned and not done (they need a TLC rerun); mechanised proof is future work \\
I8 (QUIC 0-RTT replay) is tests-only, not in TLA+ &
Model coverage of the transport-replay invariant &
Author; transport-layer implementers &
Tested through a transport-reported replay flag; vacuous where QUIC 0-RTT is not used \\
I-SIGN-JWKS is inspection-only, no test, vector or mutant &
Conformance claim for the JWKS-source invariant &
Deployment operator (inspects the bridge's key sources) &
A claim of conformance rests on inspection, not on a test \\
Key revocation is operator-only (I11-REVOKE); client cannot self-revoke &
A compromised client cannot lock the principal out; but the principal
must contact the operator to revoke a lost device key &
Deployment operator; principal &
I11-REVOKE adds operator-signed revocation, two-phase
rotation, key\_exhausted predicate; soft/hard revocation. A key\_exhausted
session keeps its timers and ends after the detachment bound unless the
operator adds a key; one whose softly revoked key stays connected is in
$a=\textsc{Att}$, has no detach timer, and ends only after the transport is
lost or the operator acts (the $a=\textsc{Res}$ case is a separate row below) \\
A re-attach signature carries no channel identifier &
A relayed genuine signature is indistinguishable from the client at the contract level &
Deployment operator; transport binding &
The client signs only nonces it requested from the intended bridge and the transport authenticates the bridge; a channel-bound signature (for example a TLS exporter value) is not specified \\
Scope is a flat set of opaque identifiers &
I3 and I10 compare sets by inclusion; a structured scope (for example a hierarchy) is outside the document unless mapped to identifiers whose subset relation is its order &
Deployment operator &
Stated in Section~\ref{sec:spec}; defining a partial order is future work \\
I14 protects only against honest rendering drift &
A malicious client that reports whatever hash it was given is not detected by I14 &
Deployment operator; client-side implementation &
The bridge has no independent view of what the client displayed; future work may bind the answer to a client-side attestation \\
Wire format for client--bridge messages specified normatively (Appendix~\ref{app:wire-format}); transport binding not specified &
Byte-level interoperability between bridges for the JSON message content; transport-layer interop &
Implementer (must choose a transport binding) &
The wire format (Appendix~\ref{app:wire-format}) specifies the wire format; a future revision MAY add a binary encoding (protobuf, CBOR) and transport bindings \\
The key registry and the event log are written in two steps &
A crash between the registry write and the commit of \texttt{ATTACHED}, \texttt{CLIENT\_ADDED} or \texttt{OPERATOR\_ADD\_KEY} leaves a registry entry the log does not record &
Storage implementer &
Recovery drops it: the log is the source of truth, so recovery repairs the registry from the log (writes back a missing entry, drops an extra one, replaces one that differs) and names the key ids in \texttt{BRIDGE\_RECOVERED}; the session is not quarantined, and a registry that has been tampered with is silently restored, so storage tampering is detected only through the log; the reference bridges undo the registration when the append fails and repair the rest at recovery; a durable store may write both in one transaction; tests T108--T112 and T143, vectors V102--V104 and V129 \\
A grant is bound to a bridge identity that the deployment configures &
Two bridges with the same identity accept each other's grants; a captured attach (grant and proof) can create the session once at its bridge, with the same key, which gives the replayer nothing it can sign with &
Deployment operator &
Give each bridge a distinct identity; the attach proof has no nonce, so freshness of the attach itself is not claimed \\
Truncation or rewriting of the last events is not detected by I5 &
A log cut back to an earlier event, or whose final event is replaced, is a valid chain; recovery accepts it, so a terminated session can become active again and a revoked key can be accepted again &
Bridge operator; storage implementer &
Hold the head (sequence and hash) outside the log, for example as an I-CHECKPOINT anchor; the reference bridges keep no external head \\
A re-attach request lets an unauthenticated caller tell a live session from a dead one &
An unknown, a terminated and a quarantined session get one refusal (\texttt{refused}), but a session that is live answers a request with a bad signature, a stale nonce or a wrong prefix under those reasons, and a key\_exhausted session says so; a caller who guesses session identifiers can therefore list the identifiers that are live, though it learns nothing about their keys (an unregistered key is answered as a bad signature) &
Transport binding; deployment operator &
Use unguessable session identifiers and rate-limit unauthenticated resume requests per source; the contract bounds only what such requests write to the log (I11g), not what they reveal \\
A join tells an attached session from a detached one before authentication &
In both reference bridges a \texttt{join} with any signature is answered \texttt{bad\_signature} (and one \texttt{RESUME\_REJECTED} is logged) when the session is attached, and \texttt{refused} with nothing logged when it is detached or re-attaching, because the cell is checked before the signature; a key\_exhausted or expired session is likewise reported before any signature check; this departs from I11 (d), which gives such answers only after the signature verifies &
Author (reference bridges); deployment operator &
Checking the signature and nonce before the cell and the timers is planned and not done; until then the mitigation of the row above applies \\
A request refused before authentication can be accepted later (pre-registration replay) &
A join or re-attach signed by a key the session does not yet register is refused (\texttt{unknown\_key}, answered \texttt{bad\_signature}; on a key\_exhausted session \texttt{key\_exhausted}, before any check) and its nonce stays live; if \texttt{CLIENT\_ADDED} or \texttt{OPERATOR\_ADD\_KEY} registers that key within $T_{\text{nonce}}$, anyone who captured the bytes can present them and is accepted. After \texttt{OPERATOR\_ADD\_KEY} on a detached, exhausted session the replayer becomes the only attached client and receives the replay of the log. Both reference bridges behave so, and the TLA+ model reaches the same trace (\texttt{NoForgery} admits it). This violates the single-use statement of I11 and the outsider statement of Section~\ref{sec:invariants-token-auth} &
Author (reference bridges, model); deployment operator until fixed &
A rule that a registry change voids the live nonces of the session, stated in I11 and checked in the model, is planned and not done. Until then: register a key before its device presents a request with it, or let $T_{\text{nonce}}$ pass after a refused attempt before registering the key \\
Replays of a captured authentic request grow the log without bound &
A replayed authentic request verifies, is recorded (\texttt{nonce\_used}) and restores the I11 (g) allowance, so a caller without a key that holds one captured request makes the log grow linearly in its effort (Appendix~\ref{app:attacker}) &
Deployment operator &
Rate-limit or authenticate the transport in front of re-attach and join; the contract has no identity to limit by before a signature \\
A caller without a key can end a detached session &
Holding the live-nonce limit (I11 (f)) refuses a detached client every challenge; kept up until $T_{\text{detach}}$ elapses, the session ends with \texttt{DETACH\_TIMEOUT} (Residual denial in Section~\ref{sec:invariants-token-auth}); L1 bounds detachment, it does not give availability &
Deployment operator &
Rate-limit or authenticate the transport in front of \texttt{challenge} \\
A later approval revives a lapsed one &
Both reference bridges keep one approval expiry per session, set by the latest \texttt{HUMAN\_RESOLVED}; a queued call approved under an escalation whose $T_{\text{approve}}$ has lapsed is forwarded at the drain that follows a later, unrelated approval, instead of \texttt{TOOL\_CALL\_EXPIRED} (\texttt{approval\_ttl}) as I13 and Table~\ref{tab:timers} require. The TLA+ model has the same single countdown, so TLC cannot see it &
Author (reference bridges, model) &
Recording the approval expiry per call at approval time, in both bridges and the model, is planned and not done; $T_{\text{call}}$ still bounds every queued call \\
The Python bridge can append \texttt{CHECKPOINT} after the end of a session &
With checkpoints enabled (the default configuration), the Python bridge emits a \texttt{CHECKPOINT} whenever the threshold is reached, with no check of the lifecycle or of quarantine and before chain verification: after \texttt{TERMINATION} (against I6 and the termination effects), on a quarantined session (against rule~Q), and over a tampered log when a re-attach starts. The suite runs with checkpoints off, so no test sees it. The Rust crate emits no checkpoints &
Author (Python bridge) &
Guarding emission by lifecycle and quarantine and placing it after chain verification, with tests that run with checkpoints on, is planned and not done; a reader of a Python bridge log ignores a \texttt{CHECKPOINT} after \texttt{TERMINATION} or after the failure event \\
A re-attach in the tick in which $T_{\text{detach}}$ expires is refused &
The edge-case rule of Section~\ref{sec:timeouts} says the request MUST be processed first; both reference bridges evaluate the timers before the request, so the session ends with \texttt{DETACH\_TIMEOUT} and \texttt{TERMINATION} and the request is answered \texttt{terminated} &
Author (reference bridges) &
Verifying the request before the timers is planned and not done; the effect is confined to the instant at which $T_{\text{detach}}$ expires \\
A soft revocation of the key being verified does not abandon the verification &
Through the two-phase interface (a re-attach begun and completed in separate calls), a \texttt{KEY\_REVOKED} with reason \texttt{rotation} between the two leaves the revoked key attached when the verification completes; if it was the only key, the session is key\_exhausted in $a=\textsc{Res}$, the completion is refused and no detach timer runs, so the session stays in $a=\textsc{Res}$ until the transport is lost, a human timer fires or the operator acts (against L1 and the bound on $a=\textsc{Res}$). The wire front-ends of both bridges run both phases in one call, so the case is not reachable over their wire; by reading, the model has the same shape (no TLC run of the property) &
Author (reference bridges, model); implementer of an asynchronous bridge &
Abandoning the verification on a soft revocation as on a hard one, in both bridges and the model, is planned and not done; until then revoke a key that is being verified with reason \texttt{compromise} \\
An operator key addition is accepted on a quarantined session &
Rule~Q permits only an operator \texttt{TERMINATION} on a quarantined session; both reference bridges accept \texttt{OPERATOR\_ADD\_KEY} with a valid operator signature on a session that is quarantined and key\_exhausted and append it. An operator-signed termination of a quarantined session that still has active keys is refused (\texttt{not\_key\_exhausted}); the unsigned \texttt{terminate} ends it &
Author (reference bridges); deployment operator &
Refusing \texttt{OPERATOR\_ADD\_KEY} on a quarantined session is planned and not done; until then the operator adds no key to a quarantined session \\
A call forwarded directly is neither scope-checked nor recorded &
While the session is in $(\textsc{Att},\textsc{None})$, both reference bridges forward a submitted call with no event and no check against the current scope, also after a narrowing or for a permission never granted; I3 and I10 bind the scope the bridge holds and the queue and drain path, not the direct path, and no MUST covers the direct path &
Runtime; deployment operator &
The runtime enforces the scope on what it executes; a MUST that a direct call outside the current scope is refused and recorded would be a design change (both bridges, vectors, model) and is not made \\
A duplicate submission leaves no record &
A call submitted for queuing (while $h=\textsc{Pend}$, or while detached) under a key the session has already queued is answered \texttt{duplicate\_key} and commits no event, so the log does not show that the runtime resubmitted the key, with the same content or other content; a key the bridge generates (\texttt{tc-} and a counter) can coincide with one the runtime supplied, and that distinct call is then dropped as a duplicate &
Runtime; author &
The runtime supplies its own key for every call or for none; generated keys that skip taken identifiers are planned and not done \\
Approvals and resolutions are not signed &
Anyone who can speak the transport can inject \texttt{resolve}, \texttt{answer\_tool\_call}, \texttt{escalate}, \texttt{disconnect} or \texttt{terminate}; the contract binds an approval to a request hash (I14), not to a person &
Transport binding; deployment operator &
These ops stay transport-authenticated (Appendix~\ref{app:wire-format}); the deployment documents the authentication and the identity accepted; the reference wire server takes a caller-identity hook but ships none \\
A key that a connected client registers is not checked for well-formedness &
\texttt{add\_client} verifies the signature of the connected key, not the scheme tag or the length of the new public key; a key of an unknown scheme or of the wrong length is registered, can never prove possession, and counts as an unrevoked key, so the session does not reach \texttt{key\_exhausted} (and the operator's key addition and termination stay unavailable) until the operator revokes it &
Client implementer; deployment operator &
The registering key is itself registered and connected; the timers and the unsigned terminate still end the session; a deployment can refuse unknown schemes in its verifier \\
Quarantine is not a cell in $Q$ (the product machine) &
The state machine has active cells and $\bot$; quarantine is a side predicate &
Reader (may miscount states) &
Stated explicitly in Section~\ref{sec:product-model} and rule Q; vectors V19, V29 exercise it \\
\bottomrule
\end{longtable}

\paragraph{Claims and their limits.} Table~\ref{tab:claims} lists each
headline claim of this document next to the evidence behind it and the
limit of that evidence.

\begin{table}[htbp]
\centering
\caption{Claims, evidence and limits.}
\label{tab:claims}
\footnotesize
\begin{tabularx}{\columnwidth}{X X X}
\toprule
\textbf{Claim} & \textbf{Evidence} & \textbf{Limit} \\
\midrule
The obligations in G1--G4 are not assigned to the bridge by AGTP, GAR, or PACE &
The comparison matrix in Section~\ref{sec:mapping}, built from the texts read
(\texttt{paper/sources/}, with \texttt{paper/sources/MANIFEST.md}
recording the URL and SHA-256 of each version; the sources are working
copies and are not re-licensed under CC-BY-4.0 or MIT) &
An argument from the drafts as read; it weakens if one of them is extended.
Cells not verified are marked as such in the matrix. \\
The product state machine and invariants I1--I7 and I9--I15 have no violation under bounded model checking, with I8 not modelled and I6 and I11 (f) only weakly checked (see the register) &
TLC on bounded configurations; mutants of the TLA+ model are rejected
(Appendix~\ref{app:formal}) &
Bounded sessions, clients and permissions; hashes and event
replay are abstracted, so I4, I5 and I14 are checked as flags only, I8 is tests-only, the state invariant for I6 is a tautology and the no-eviction property of I11 (f) is vacuous as configured, and signatures
are a possession predicate with unforgeability assumed, and the attacker is checked only in a narrow configuration (one permission, one call key, no log tampering, two remembered signatures); no mechanised proof for unbounded parameters;
liveness only under stated fairness. \\
The invariants are testable &
Pytest suite T1--T158, property-based tests T8--T10, a killed mutant for
every invariant I1--I15 (Table~\ref{tab:invariant-independence}) &
Tests run against the author's own bridge; a mutant kill shows the suite
is sensitive to that edit, not that the invariant holds in general. \\
One author implemented the specification from the text &
Rust crate written by the same author after the Python bridge, from
Sections 4 and 5 and with comments, error texts and the generated key
format shared with the Python code, so it is not an independent reading; shared vectors V01--V134 pass on the Python bridge, and on the
crate except V52--V54, V88 and V89, which expect \texttt{CHECKPOINT}
events the crate does not emit and which its runner skips; the crate runs the same vectors on two storage backends &
One author for both implementations; agreement shows one reading is
consistent, not that the text is unambiguous. A second, independent
implementation is the next step; V01--V134 are the interface for it. \\
Proof of possession (I11) holds against an outsider, except for the
pre-registration replay in the register; I11-REVOKE
extends it to revoked keys &
Ed25519 and ECDSA P-256 (wrong key, other message, key added after the
challenge; the proof of the other operation, another key id or no tag;
V59--V62, V90--V94, T82--T84, T101--T104) in the Python bridge and the Rust crate;
key identifiers scoped to a session and a registry repaired from the log at recovery (T105--T112, T143, V95, V96, V102--V104, V129, mutants \texttt{registry-*}); a proof of possession at attach (T113--T118, V97, V98, V101);
registering a further key needs a signature by a connected, unrevoked key over the key, the session and a single-use nonce (unsigned, forged, other-operation, replayed, expired and revoked-key requests are refused with a named reason: V105--V108, T119--T123, mutants \texttt{I15-add-client-*});
registry holds public keys and a scheme tag only; mutants for forged requests; revocation, rotation, key exhaustion and in-flight revocation: tests T67--T76, T85--T89 and T134--T137, vectors V63--V74 and V115--V120 on both implementations (the timers of a key\_exhausted session keep running: mutants \texttt{TIMERS-key-exhausted-*}, \texttt{EXHAUST-*}); the operator signature on a revocation, key addition and termination is verified (T90--T96, V75--V85, mutants \texttt{IOPERATOR-*}); an unsigned terminate is refused on a key\_exhausted session (T145, T146, V130, mutants \texttt{I11-REVOKE-bare-terminate-*}, \texttt{EXHAUST-bare-terminate-*}), and a re-attach for an unregistered key is answered as a bad signature (T147, V131, mutants \texttt{I11-unknown-key-distinguishable}, \texttt{RESUME-unknown-key-distinguishable});
two sessions on one bridge do not share the resources listed (nonces, challenge, rejection and revocation-request budgets, tool call keys, the authentication hook of the unsigned wire ops: T152--T155, V133, V134, mutants \texttt{session-*}, \texttt{SESSION-*}); the TLC model is single-session and does not see this class &
Unforgeability of the signature scheme is assumed; the model's attacker holds no
client key (bounded as in the register). It does not hold for a request
refused before its key was registered: both bridges and the model accept
its replay once the key is registered within $T_{\text{nonce}}$ (register). The operator key is the one that signs scope grants; its
custody, rotation and revocation are the deployment's, and a signature
carries no sequence number, so a captured one is replayable only where
the operation is still legal (a key already revoked or registered is
refused). \\
The initial scope is verified against a grant before attach is committed &
Tests T77--T81 and T113--T118, vectors V55--V58 and V99--V101 (the grant binds the public key and the bridge) and mutants of the check
(Table~\ref{tab:invariant-independence}) &
The grant format is the deployment's; the reference verifier checks an
operator signature, so what a real authority issues is not tested; the
model does not represent attach. \\
The contract is usable in a product runtime &
Mapping onto Termenbox, which changed the contract &
Same author; the server is not wired to the crate; no production case
study. \\
\bottomrule
\end{tabularx}
\end{table}

\paragraph{Threats to validity.} The safety argument in
Section~\ref{sec:safety-argument} is informal. TLC checks the invariants
only on bounded models, and the cryptographic checks are not modelled
there, so the unforgeability of the signature scheme is an assumption
and is not model-checked; the implementations exercise it with Ed25519
against a registry of public keys (Appendix~\ref{app:formal-limitations}). The conformance suite,
its mutants and the vectors test two implementations (the Python
reference bridge and the Rust crate); they show consistency of one reading
of the specification, not independent confirmation, and no third-party
bridge has been tested. The specification
does not address cross-session correlation (deferred to
AGTP-SESSION \S 10.3) or governance of multi-agent composition
(deferred to A2A and MCP composition analyses). The product machine
(Section~\ref{sec:product-model}) is the normative state model; the
state names that survive in the figures are a lossy derived view,
not the semantics. No case study from a production runtime exists yet,
and the author's runtime project (Termenbox) intends to adopt this
specification, without a schedule, so the first production case study is still to come.

% ================================================================================
